\begin{textsample}{NEG Dim 2 – ba – Score 1.00 – ba/ba\_inf\_6.txt}  \label{ex:f2_neg_260}
7.6.
Entretecimentos entre Currículo, Didática e Processos de Aprendizagem na Educação Infantil Atos de currículo, mediação didática e processos de aprendizagem são entretecimentos do campo educacional que guardam relações profundamente necessárias.
Podemos dizer que currículos, para se qualificarem através de seus atos, precisam de mediações didáticas qualificadas, condição para a possibilidade de aprendizagens e processos formativos também qualificados.
No caso da Educação Infantil, currículos e seus atos precisam que a didática assuma, de início, as especificidades das aprendizagens e do desenvolvimento infantis, bem como as transversalidades que fundam a educação das crianças, como a concepção de infância, o cuidado, a ludicidade, a interação acolhedora e motivacional, para que o ensino não perca a compreensão de que lida via mediações com as infâncias em educação.
Métodos ativos e construcionistas que tomem as crianças como autoras das suas aprendizagens, que compreendam a aprendizagem como um fenômeno ao mesmo tempo cognitivo, social, cultural, ético e político podem qualificar os atos de currículo por meio dos entretecimentos com a didática e a compreensão do processo de aprendizagem das crianças.
Educadores atenciosos, acolhedores, brincantes, mediadores de aprendizagens criativas, facilitadores das experiências aprendentes das crianças, conectados com uma Educação Infantil contemporânea e cientes da pluralidade, dos espaços-tempos de formação dos nossos tempos, podem trazer para a Educação Infantil entretecimentos ricos e valorosos entre currículo, didática e processos de aprendizagem.

% matched lemmas: 
\end{textsample}
